\section{A target-relative reduction in the fixed-width cold limit}
\label{sec:cold-limit-reduction}

The unrestricted Gaussian question can be reduced to an approximation
energy comparison without interchanging width, depth, and variance limits.
This reduction does not by itself provide the missing function-level
energy lower bound.

Fix dimension $d$, width $N$, and depth $L$. In the energy coordinates of
the model the raw Gaussian is $\mathcal N(0,\sigma^2 I)$, and the squared
Euclidean norm is $\kappa_N$. For a Boolean function $g$, define
\[
 e_{N,L}(g)=\inf\{\kappa_N(w):g(x)z_w(x)\ge1/2
                           \text{ for all }x\in\{-1,1\}^d\},
\]
with infimum $+\infty$ if the set is empty.
These are approximation costs of Boolean functions, not costs of specified
implementations.

\begin{proposition}[Exact fixed-width cold posterior rates]
\label{prop:cold-posterior-rates}
Let $S$ be any fixed noiseless labeled sample for which the posterior is
defined, and let $\mathcal V_S$ be the Boolean functions consistent with
its labels. Put $e_S=\min_{g\in\mathcal V_S}e_{N,L}(g)$.
For every subset $\mathcal A$ of Boolean functions with
$\mathcal A\cap\mathcal V_S\ne\varnothing$,
\[
 \lim_{\sigma\downarrow0}-2\sigma^2\log\Pi_{N,\sigma}
               (\mathcal A\mid S)
 =\min_{g\in\mathcal A\cap\mathcal V_S}e_{N,L}(g)-e_S.
\]
The identity is interpreted in the extended real sense when all functions
in the numerator have infinite cost. In particular, the cold posterior
concentrates on the functions of minimum energy in the version space.
At a fixed $N\ge2$, this cold limit is also obtained along the exponent
convention $\sigma^2=N^{-\gamma}$ by taking $\gamma\to\infty$.
\end{proposition}

\begin{proof}
The finite-dimensional Gaussian large-deviation principle has speed
$1/(2\sigma^2)$ and rate $\kappa_N$. The strict-margin event
$A_g=\{g(x)z_w(x)>1/2\text{ for all }x\}$ is open. Its closure is
exactly the corresponding closed-margin event: for a network in the
latter, multiplying its readout by $1+\eta$ puts it in $A_g$, and its
energy tends to the old energy as $\eta\downarrow0$.
Thus the upper and lower large-deviation bounds agree, giving
$-2\sigma^2\log\mu_{N,\sigma}(A_g)\to e_{N,L}(g)$.
The globally margin-conditioned, sign-consistent posterior has function
probabilities proportional to $\mu_{N,\sigma}(A_g)$ for
$g\in\mathcal V_S$. There are finitely many Boolean functions at fixed
$d$, so each sum has rate equal to its least constituent energy. Taking
the ratio proves the claim. Coercivity and closedness ensure that every
finite minimum energy is attained.
\end{proof}

\begin{lemma}[A uniformly cheap sample memorizer]
\label{lem:sample-memorizer-energy}
For integers $d,n\ge1$, set
\[
 C_n=2n+1,\quad b=\log(16(n+1)),\quad R_n=4+(d+1)b,
\]
\[
 M=\left\lceil\max\{8,C_n,2R_n,dC_n(d+1)R_n^2\}\right\rceil,
 \qquad T=20M(n+1).
\]
Every noiseless labeled sample of size at most $n$ has a consistent
Boolean memorizer with unit-margin realizations at every width $N\ge M$
and every depth $L\ge2$, of energy at most
\[
 K_{\rm mem}(d,n,L)=
 \begin{cases}
  2+2T,&L=2,\\
  2+(L+14)(2T)^{1/(L-1)},&L\ge3.
 \end{cases}
\]
Here $M,T$ are polynomial in $d,n$. In particular, if
$L\ge3$ and $L-1\ge\log(2T)$, then
\[
 K_{\rm mem}(d,n,L)\le L+B(d,n),\qquad
 B(d,n)=16+32\log(2T)=O(\log(d+n)).
\]
\end{lemma}

\begin{proof}
Let $P$ be the set of distinct positively labeled sample inputs, and put
$p=|P|\le n$. For each $a\in P$, set
\[
 h_a(x)=\tanh\bigl(b(a\cdot x-d+1/2)\bigr).
\]
Its distance from $2\mathbf1\{x=a\}-1$ is at most
$2e^{-b}=1/[8(n+1)]$. Consequently
\[
 z(x)=2\left[-1+\sum_{a\in P}(h_a(x)+1)\right]
\]
has signed margin at least $7/4$ for the Boolean function that is positive
on $P$ and negative elsewhere. It fits all sample labels.
Implement the constant term using $p+1$ constant hidden neurons, each
with activation $1/2$ and readout weight $4(p-1)/(p+1)$.
The $p$ isolator readout weights are two. Thus there are $2p+1\le C_n$
neurons, all weights are bounded by $R_n$, and
\[
 \|v\|_2^2=4p+16(p-1)^2/(p+1)\le20(n+1).
\]
The construction includes $p=0$, where it is the constant negative
classifier. Zero-padding to width $M$ gives first-layer energy at most
one and readout energy at most $T$. Balanced cloning to every $N\ge M$
increases each by at most a factor of two. The shallow scalar-transmission
construction, with its final readout doubled, therefore gives the stated
$K_{\rm mem}$ at every depth. Its constants and argument are exactly those
in the shallow-target theorem from the preceding round.
Finally, putting $u=\log(2T)/(L-1)\le1$ and using $e^u\le1+2u$ gives
\[
 K_{\rm mem}\le L+16+
       2\frac{L+14}{L-1}\log(2T)
 \le L+16+32\log(2T).
\]
\end{proof}

For a test distribution $D$, target $f$, and $0<\varepsilon<1$, define
the best energy of an accurate Boolean function by
\[
 e_{N,L}^{\varepsilon}(f;D)
 =\min_{g:\Pr_D(g(X)\ne f(X))\le\varepsilon}e_{N,L}(g).
\]

\begin{corollary}[A sufficient energy barrier for posterior failure]
\label{cor:cold-approximation-barrier}
Fix $d,n,L$, a width $N\ge\max\{2,M,C_L(f)\}$, a target $f$, and a
test distribution $D$. If for some $\Delta>0$,
\[
 e_{N,L}^{\varepsilon}(f;D)
     \ge K_{\rm mem}(d,n,L)+\Delta,
\]
then, uniformly over all size-$n$ noiseless training samples from $f$,
the posterior mass of functions with error at most $\varepsilon$ tends
to zero as $\sigma\downarrow0$. More precisely, for all sufficiently
small $\sigma$ this mass is at most $\exp[-\Delta/(4\sigma^2)]$.
The posterior Gibbs error consequently has limit inferior at least
$\varepsilon$. The conclusion holds along arbitrarily large
$\gamma$ under $\sigma^2=N^{-\gamma}$.
\end{corollary}

\begin{proof}
The memorizer ensures $e_S\le K_{\rm mem}$ for every sample. The exact
cold posterior rate of the accurate set is therefore at least $\Delta$.
Its mass is eventually at most the displayed exponential. There are only
finitely many input sequences of length $n$ at fixed $d$, so the cold
threshold can be chosen uniformly over all such samples. Thus it is not
a sample-dependent choice of prior scale. The Gibbs-error conclusion
follows by assigning error at least $\varepsilon$ to the complement.
\end{proof}

This is a genuine conditional posterior comparison, but its hypothesis
is not established for a small-circuit target. A large cost for one
particular implementation, a cloning obstruction, or a lower bound only
for the exact target $f$ is insufficient: the barrier must cover every
accurate competing Boolean function, with arbitrary allowed weights at
the chosen width and depth. For depths exceeding a constant multiple of
$\log(d+n)$, the missing barrier is above an energy benchmark
$L+O(\log(d+n))$. An explicit small-circuit family satisfying such an
approximation barrier would resolve a substantive part of the unrestricted
question negatively. No such family is proved here.
